\chapter{GIỚI THIỆU}

\section{Tổng quan vấn đề}

Trong bối cảnh kinh tế toàn cầu hóa, ngành thẻ tín dụng đang phát triển với tốc độ chóng mặt. Các ngân hàng và tổ chức tài chính phải đối mặt với thách thức cấp thiết để hiểu rõ hơn về khách hàng của họ để có thể:

\begin{itemize}
    \item Phát triển các sản phẩm dịch vụ phù hợp
    \item Tối ưu hóa chiến lược tiếp thị
    \item Quản lý rủi ro tín dụng hiệu quả
    \item Cải thiện tỷ lệ giữ chân khách hàng
    \item Tăng thu nhập từ các dịch vụ tài chính
\end{itemize}

Phân cụm khách hàng (Customer Segmentation) là một kỹ thuật quan trọng trong khai phá dữ liệu cho phép các tổ chức nhận diện và hiểu các nhóm khách hàng khác nhau để áp dụng các chiến lược khác nhau.

\section{Mục tiêu của báo cáo}

\textbf{Mục tiêu chính:} Phân tích, phân cụm và đặc trưng hóa các khách hàng thẻ tín dụng để hỗ trợ các quyết định kinh doanh và phát triển chiến lược tiếp thị được cá nhân hóa.

\textbf{Các mục tiêu cụ thể:}

\begin{enumerate}
    \item \textbf{Xử lý dữ liệu:} Làm sạch, chuẩn hóa và chuẩn bị dữ liệu thẻ tín dụng cho phân tích
    
    \item \textbf{Khám phá dữ liệu:} Thực hiện phân tích khám phá toàn diện để hiểu đặc điểm của dữ liệu, các phân phối, mối tương quan
    
    \item \textbf{Xác định số lượng cụm tối ưu:} Sử dụng các phương pháp như Elbow Method và Silhouette Analysis để xác định K tối ưu
    
    \item \textbf{Phân cụm khách hàng:} Áp dụng K-means clustering để phân khúc khách hàng thành các nhóm đồng nhất
    
    \item \textbf{Xác thực kết quả:} So sánh kết quả với các phương pháp phân cụm khác (phân cụm phân cấp, DBSCAN)
    
    \item \textbf{Phân tích đặc trưng:} Phân tích chi tiết các đặc trưng của từng cụm
    
    \item \textbf{Phát hiện dị biệt:} Xác định các khách hàng dị biệt có hành vi khác thường
    
    \item \textbf{Đưa ra khuyến nghị:} Đề xuất các chiến lược tiếp thị và kinh doanh dựa trên phân cụm
\end{enumerate}

\section{Phạm vi của báo cáo}

Báo cáo tập trung vào:

\begin{itemize}
    \item \textbf{Dữ liệu:} Bộ dữ liệu công khai về thẻ tín dụng gồm 8,950 khách hàng và 17 biến tài chính
    \item \textbf{Phương pháp:} Các kỹ thuật khai phá dữ liệu không giám sát (unsupervised learning)
    \item \textbf{Công cụ:} Python với các thư viện scikit-learn, pandas, matplotlib, seaborn
    \item \textbf{Khoảng thời gian:} Phân tích dữ liệu tổng hợp (snapshot) từ một thời điểm
\end{itemize}

\section{Cấu trúc báo cáo}

\begin{itemize}
    \item \textbf{Chương 1 - Giới thiệu:} Tổng quan bài toán, mục tiêu và phạm vi
    \item \textbf{Chương 2 - Tổng quan tài liệu:} Nền tảng lý thuyết cho phân cụm, các phương pháp chính
    \item \textbf{Chương 3 - Phương pháp:} Các kỹ thuật và thuật toán được sử dụng
    \item \textbf{Chương 4 - Xử lý dữ liệu:} Làm sạch dữ liệu, chuẩn hóa, biến đổi
    \item \textbf{Chương 5 - Phân tích khám phá:} Thống kê mô tả, trực quan hóa, phân tích mối tương quan
    \item \textbf{Chương 6 - Kết quả phân cụm:} Áp dụng K-means, xác định K tối ưu, phân tích cụm
    \item \textbf{Chương 7 - Thông tin chi tiết kinh doanh:} Đặc trưng của các cụm, phát hiện dị biệt, khuyến nghị
    \item \textbf{Chương 8 - Kết luận:} Tóm tắt kết quả, hạn chế, hướng phát triển tương lai
    \item \textbf{Phụ lục:} Mã nguồn, bảng dữ liệu chi tiết, đặc tả kỹ thuật
\end{itemize}

\section{Những đóng góp chính}

\begin{itemize}
    \item Cung cấp một phương pháp toàn diện cho phân cụm khách hàng thẻ tín dụng
    \item Xác định được các phân khúc khách hàng có ý nghĩa kinh doanh
    \item Đưa ra các chiến lược tiếp thị được cá nhân hóa dựa trên phân cụm
    \item Phát hiện các khách hàng dị biệt có tiềm năng hoặc rủi ro cao
    \item Cung cấp một mô hình có thể được dùng để phân loại các khách hàng mới
\end{itemize}
